\subsection{Inherent variability in biological systems}

\paragraph{Overview: sources of variability in biology}

In all the sections described before, we applied the same biological model to the patients of a given cohort. However, we observe strong variability between the samples, even within the same individual. The variability occurs at three levels: at the environmental level (disease state, tissue location, . . . ), the genotypical level (presence of individual mutations inducing varying transcriptomic activity) and even at the cell population level. 

In the next sections, we will first detail canonical clustering methods, namely \Glspl{gmm} that are still relevant in the exploratory stage to unravel homogeneous groups in a seemingly highly variable disease condition, both at the phenotypical and transcriptomic level. Then, we review some of the most common deconvolution methods designed to measure and for some of them, address the bias induced by the strong variability of the cell composition.



\subsection{Mixture models} 
\paragraph{Overview: introduction to mixture models}

\subsubsection{Gaussian mixture models}

\subsubsection{The EM algorithm}
\paragraph{Motivation}
\paragraph{Derivation of the EM algorithm in \glsfmtshort{gmm}}
\paragraph{Clustering in high-dimensional setting}

\subsubsection{Statistical results applied to the EM algorithm}

\paragraph{Model selection}
\paragraph{Confidence intervals}
\paragraph{Extensions of the EM algorithm}



\subsection{Cellular deconvolution}
\paragraph{Overview: cellular deconvolution, a promising but challenging method to unravel complex biological structures}

\subsubsection{Physical approaches}

\subsubsection{Computational approaches}

\paragraph{Reference-free approaches}

\paragraph{Marker-based approaches}

\paragraph{Partial-based approaches}



